\chapter{The Macroscopic Oscillator and the Structural Phase Limits}
\label{ch:macroscopic_oscillator}

Standard cosmology describes large-scale cosmic evolution through the dynamical FLRW metric and, within the standard $\Lambda$CDM framework, extrapolates this evolution toward progressively dilute future states. Tension Field Theory (TFT) instead proposes that the active cosmic domain is a structurally bounded visco-elastic medium embedded within the underlying substrate and operating through recurrent thermodynamic and structural phase cycles.

In this picture, the active universe is not an indefinitely expanding geometric volume. It is a finite macroscopic State I reservoir whose structural evolution is governed by phase conversion, thermodynamic relaxation, global percolation, and eventual Catastrophic Decompression.

\section{The Optical Horizon vs. the Structural Reservoir}
\label{sec:optical_vs_structural}

The macroscopic architecture of the universe is governed by two distinct, functionally separate scales: the Viscous Cosmic Horizon, which characterizes optical relaxation, and the Structural Reservoir, which defines the mechanical limit of the active cosmic domain.

\subsection{The Viscous Cosmic Horizon ($R_{\mathrm{U}}$)}
\label{subsec:viscous_horizon_macro}

As established in \hyperref[chap:cosmological_anomalies]{Chapter~\ref*{chap:cosmological_anomalies}}, the characteristic cosmological relaxation scale
\begin{equation}
    R_{\mathrm{U}}
    =
    \frac{c_{\mathrm{E}}}{\Gamma_{\mathrm{D}}}
    \approx
    1.267 \times 10^{26}
    \text{ m}
\end{equation}
defines the Viscous Cosmic Horizon.

It is the characteristic one-$e$-fold optical relaxation length of the active FCC vacuum. Visco-elastic damping continuously attenuates the coherent transverse energy of a propagating photon according to
\begin{equation}
    E_{\mathrm{obs}}(x)
    =
    E_{\mathrm{emit}}
    \exp\left(
        -\frac{\Gamma_{\mathrm{D}}x}{c_{\mathrm{E}}}
    \right),
\end{equation}
so that at the Viscous Cosmic Horizon,
\begin{equation}
    E_{\mathrm{obs}}(R_{\mathrm{U}})
    =
    \frac{E_{\mathrm{emit}}}{e}.
\end{equation}

The significance of this scale does not arise because expanding space exceeds the speed of light, nor because the physical lattice terminates there. $R_{\mathrm{U}}$ marks the characteristic distance over which optical coherence is substantially relaxed by the dissipative medium. A coherent transverse excitation may continue well beyond one $R_{\mathrm{U}}$, progressively approaching its frequency-dependent assimilation threshold into the ambient thermal reservoir represented by the Cosmic Microwave Background (CMB).

The observable cosmos is therefore not the full mechanical extent of the universe. It is a localized optical envelope within a much larger dissipative medium.

Furthermore, this macroscopic relaxation scale is directly linked to the microscopic quantum hierarchy. As established by the Scale-Invariant Continuum relation in \hyperref[sec:scale_hierarchy]{Section~\ref*{sec:scale_hierarchy}},
\begin{equation}
    E_{\mathrm{max}}
    =
    h\Gamma_{\mathrm{D}}N_{\mathrm{D}},
\end{equation}
the maximum microscopic excitation capacity and the macroscopic relaxation rate participate in the same invariant scale structure of the active vacuum. The microscopic energy ceiling and the cosmological relaxation scale therefore emerge as complementary limits of the same continuum.

Thus, $R_{\mathrm{U}}$ is not the strict spatial boundary of State I. It is the characteristic kinematic relaxation scale of the active state.

\subsection{The Structural Reservoir ($R_{\mathrm{Res}}$)}
\label{subsec:structural_reservoir_macro}

However, the active FCC lattice does not terminate where ordinary optical access begins to fade.

As derived through the Vacuum Phase Conversion mechanism in \hyperref[part:two]{Part II}, the total volumetric capacity required to structurally support the active cosmic domain reveals a vastly larger underlying lattice reservoir. The FCC structure extends to an absolute mechanical macro-limit defined as the Structural Reservoir,
\begin{equation}
    R_{\mathrm{Res}}
    \sim
    10^{35}
    \text{ m}.
\end{equation}

Rather than an arbitrary coordinate cutoff, $R_{\mathrm{Res}}$ is the mathematical boundary obtained from the mass-conservation relation governing the TFT vacuum phase transition. It is the macroscopic volume required by the phase-conversion construction to accommodate the total structural mass-energy while maintaining the active FCC phase below its saturation limit.

This produces a sharp separation between what can be seen and what must mechanically exist.

The universe is not expanding through empty geometry; it is a structurally contiguous tension field whose macroscopic evolution is mechanical rather than metric. Its active state may contract, accumulate phase defects, reach global percolation, and undergo Catastrophic Decompression, but none of these processes requires an expanding spatial metric.

The observable cosmos is therefore a localized optical envelope embedded within a vast, stabilizing mechanical reservoir extending far beyond our optical reach.

\section{The Global-Local Energy Equivalence}
\label{sec:global_local_energy_equivalence}

The scale-invariant symmetry of the active vacuum is not merely geometric; it is also thermodynamic. This is demonstrated by comparing the integrated mass-energy of the macroscopic Structural Reservoir with the local structural yield scale of its microscopic lattice.

\subsection{The Global Cosmic Energy ($E_{\mathrm{Res}}$)}

Based on the cosmic reservoir mass derived in \hyperref[sec:reservoir_mass]{Section~\ref*{sec:reservoir_mass}},
\begin{equation}
    M_{\mathrm{Res}}
    \approx
    1.6707 \times 10^{79}
    \text{ kg},
\end{equation}
the total integrated mass-energy of the reservoir is
\begin{equation}
    E_{\mathrm{Res}}
    =
    M_{\mathrm{Res}}c_{\mathrm{E}}^2
    \approx
    1.5015 \times 10^{96}
    \text{ J}.
\end{equation}

\subsection{The Local Yield Tension ($\tau_{\mathrm{D,bulk}}$)}

We compare this global value directly with the Gravitational Yield Tension ($\tau_{\mathrm{D,bulk}}$) of the active vacuum. This quantity represents the theoretical maximum energy density supported by the active lattice at its local structural yield scale. As established in \hyperref[app:phase_transition]{Appendix~\ref*{app:phase_transition}}, the corresponding critical density,
\begin{equation}
    \rho_{\mathrm{G}}
    \approx
    2.4259 \times 10^{79}
    \text{ kg/m}^3,
\end{equation}
defines the local structural yield boundary of the FCC phase. Therefore,
\begin{equation}
    \tau_{\mathrm{D,bulk}}
    =
    \rho_{\mathrm{G}}c_{\mathrm{E}}^2
    \approx
    2.1803 \times 10^{96}
    \text{ J/m}^3.
\end{equation}

The physically meaningful cross-scale correspondence is obtained by dividing the total reservoir energy by this local yield-energy density:
\begin{equation}
    V_{\mathrm{eq}}
    \equiv
    \frac{E_{\mathrm{Res}}}
    {\tau_{\mathrm{D,bulk}}}
    =
    \frac{M_{\mathrm{Res}}}
    {\rho_{\mathrm{G}}}
    \approx
    0.6886
    \text{ m}^3.
\end{equation}

The entire accumulated mass-energy distributed across the $10^{35}$ m Structural Reservoir is physically equivalent to the yield-scale structural energy represented by only about one cubic meter of active vacuum. More precisely, approximately $0.6886$ m$^3$ of active vacuum evaluated at the critical yield density $\rho_{\mathrm{G}}$ carries the same mass-equivalent energy as the entire integrated mass-energy assigned to the macroscopic Structural Reservoir.

Within TFT, this cross-scale equivalence reveals how little of the active vacuum's enormous structural capacity is required to construct the observable universe of matter, gravitational gradients, and persistent causal structure. The macroscopic causal-information network is therefore an exceedingly sparse structural modulation of the active vacuum: localized mass and phase defects recruit only a minute fraction of the available vacuum capacity into persistent tension gradients, and those gradients organize the gravitational relationships through which spatially separated events become mechanically connected and distinguishable across $R_{\mathrm{Res}}$.

What appears macroscopically as the total mass-energy and causal structure of the universe is therefore a remarkably dilute pattern written into the vastly greater energetic capacity of the active vacuum.

\section{The Thermodynamic Ledger: Closing the Thermodynamic Cycle}
\label{sec:thermodynamic_ledger}

Without a mechanism to recover the active degrees of freedom consumed by continuous relaxation, the visco-elastic degradation of coherent transverse waves would drive the cosmic medium progressively toward thermodynamic equilibration. In TFT, the energy of a photon does not disappear as its coherent transverse state relaxes. Rather, the ordered carrier progressively loses coherent energy to noncoherent degrees of freedom of the active vacuum.

Within the present TFT CMB-reservoir interpretation, the terminal destination of this degraded optical energy is the ambient background of longitudinal lattice excitations: the phonon noise floor identified with the Cosmic Microwave Background (CMB). In this mechanical sense, the CMB acts as the long-lived ``phonon graveyard'' of exhausted transverse radiation. The microscopic derivation of complete thermalization and of the observed $2.725$ K spectrum remains an open first-principles bridge, but the thermodynamic role is clear: coherent optical structure is progressively transferred into an increasingly isotropic reservoir of noncoherent lattice excitation.

If radiative degradation ended only in the CMB phonon reservoir, the universe would still evolve toward terminal thermal equilibration. TFT therefore contains a separate structural pathway. Localized matter does not remain indefinitely embedded in the active State I medium; through its gravitational loading of the surrounding FCC lattice, it progressively nucleates Black Hole phase conversion. The CMB relaxation channel and the Black Hole phase-conversion channel are mechanically distinct, but both participate in the evolution of the same State I cosmic cycle.

The global-local energy equivalence established in \hyperref[sec:global_local_energy_equivalence]{Section~\ref*{sec:global_local_energy_equivalence}} makes this mechanism especially significant. The integrated mass-energy assigned to the entire $10^{35}$ m Structural Reservoir is equivalent to the yield-scale structural energy represented by only
\begin{equation}
    V_{\mathrm{eq}}
    =
    \frac{E_{\mathrm{Res}}}{\tau_{\mathrm{D,bulk}}}
    =
    \frac{M_{\mathrm{Res}}}{\rho_{\mathrm{G}}}
    \approx
    0.6886
    \text{ m}^3
\end{equation}
of active vacuum evaluated at its critical yield density.

Black Hole formation therefore cannot be understood in TFT as matter somehow supplying the enormous structural energy of the resulting State II domain. Matter is the trigger, not the bulk energetic substance of the phase transition.

Localized mass produces a persistent tension gradient and progressively loads the surrounding active FCC vacuum. When that local loading reaches the gravitational saturation condition,
\begin{equation}
    \alpha_{\mathrm{G}}
    \rightarrow
    1,
\end{equation}
the active FCC configuration can no longer sustain the State I geometry and undergoes the phase transition
\begin{equation}
    \mathrm{FCC}
    \rightarrow
    \mathrm{SC}.
\end{equation}

The localized matter defect therefore acts as the ignition condition for a self-sustaining crystallographic chain reaction. Once nucleated, the advancing FCC/SC phase boundary exposes additional active vacuum to conversion as the State II region grows. In TFT, a Black Hole is consequently not fundamentally matter compressed ever more deeply into itself. It is a vacuum phase-conversion process nucleated by localized mass: the matter defect provides the trigger, while the active vacuum provides the phase-convertible medium.

These localized phase collapses become the ultimate structural sinks of the macroscopic cycle. Matter defects are incorporated into the growing phase-boundary system, inertial mass becomes associated with the live FCC/SC boundary, and the enclosed region is converted into the frozen, phase-locked, zero-void State II geometry characterized by $\rho_{\mathrm{S}}$.

The thermodynamic and structural ledgers therefore advance through distinct but concurrent channels. Coherent transverse radiation progressively relaxes toward the CMB phonon reservoir, while localized mass loading independently drives FCC$\rightarrow$SC phase conversion and the growth of State II domains. The two channels do not require direct energy transfer between one another; their connection emerges at the level of the global cosmic cycle.

Crucially, the consumption of State I through phase conversion does not continue indefinitely. As Black Holes grow, their aggregate live-boundary mass,
\begin{equation}
    M_{\mathrm{p}},
\end{equation}
increases. Exterior strain fields begin to overlap, saturated FCC/SC bridges form between previously separated State II domains, and the distributed boundary system advances toward the global percolation threshold.

When
\begin{equation}
    M_{\mathrm{p}}
    \rightarrow
    M_{\mathrm{trg}},
\end{equation}
the connected live-boundary and saturated-bridge network becomes a reservoir-spanning fracture structure. The global containment differential then collapses,
\begin{equation}
    \Delta P
    \rightarrow
    0,
\end{equation}
triggering the Catastrophic Decompression---the Big Rupture---detailed in \hyperref[chap:the_big_rupture]{Chapter~\ref*{chap:the_big_rupture}}.

Within the TFT cyclic closure, terminal Heat Death is therefore mechanically preempted. The universe does not escape thermodynamic degradation by locally reversing entropy. Rather, radiative degradation accumulates within the CMB phonon reservoir while an independent mass-triggered phase-conversion process progressively converts the active lattice into a globally unstable State II network. Once that network reaches structural percolation, the configuration in which the thermodynamic degradation accumulated can no longer be mechanically sustained.

The complete thermodynamic and structural sequence may therefore be summarized as
\begin{equation}
    \begin{aligned}
        &\text{coherent transverse radiation}
        \longrightarrow
        \text{CMB phonon reservoir},
        \\
        &\text{localized mass loading}
        \longrightarrow
        \text{State II phase conversion}
        \longrightarrow
        \text{global percolation}
        \longrightarrow
        \text{Catastrophic Decompression}
        \longrightarrow
        \text{renewed State I}.
    \end{aligned}
\end{equation}

The universe consequently operates as a macroscopic oscillator in which the thermodynamic and structural clocks advance together without being the same process: light ultimately relaxes into the phonon graveyard, localized mass ignites the State II chain reaction, the growing State II network develops into a global fault structure, and that fault structure ultimately forces the mechanical renewal of the active vacuum.

That is the closure of the TFT thermodynamic ledger.

\section{The Chronology of the Macroscopic Oscillator}
\label{sec:chronology_oscillator}

Within standard $\Lambda$CDM, the approximately $13.8$ billion-year cosmic age follows from the integrated cosmological expansion history. TFT assigns a different role to the nearby but distinct timescale obtained from the locally measured Hubble slope. Under the present empirical calibration,
\begin{equation}
    \Gamma_{\mathrm{D}}
    =
    H_0^{(\mathrm{obs})},
\end{equation}
so that
\begin{equation}
    \frac{1}{\Gamma_{\mathrm{D}}}
    \approx
    13.387
    \text{ billion years}.
\end{equation}

In TFT, this quantity is the characteristic vacuum relaxation time, not the age of the active cosmic reservoir. The much longer structural chronology follows instead from the accumulated FCC$\rightarrow$SC phase conversion required to reach global percolation.

\subsection{The Global Phase Avalanche Equation}
\label{sec:phase_avalanche_equation}

Black Hole formation in TFT is a nucleated phase transition. Once a local State II domain is formed, continued growth proceeds through the self-sustaining advancement of the FCC$\rightarrow$SC phase boundary.

The chronology adopts the single-rate macroscopic relaxation closure developed in \hyperref[app:energy_and_chronology]{Appendix~\ref*{app:energy_and_chronology}},
\begin{equation}
    \frac{1}{M_{\mathrm{p}}}
    \frac{\mathrm{d}M_{\mathrm{p}}}{\mathrm{d}t}
    =
    \Gamma_{\mathrm{D}}.
\end{equation}

Integration gives the Global Phase Avalanche Equation,
\begin{equation}
    M_{\mathrm{p}}(t)
    =
    M_{\mathrm{p}}(0)
    e^{\Gamma_{\mathrm{D}}t},
\end{equation}
where $M_{\mathrm{p}}(0)$ is the primordial aggregate live-boundary mass established during the short State II nucleation interval following Catastrophic Decompression.

The equality between the phase-avalanche fractional growth rate and $\Gamma_{\mathrm{D}}$ is the single-rate macroscopic closure of the present TFT chronology. It does not imply direct conversion of coherently relaxed photon energy into Black Hole mass; the optical-relaxation and phase-conversion channels remain mechanically distinct.

\subsection{The Total Oscillation Time ($T_{\mathrm{cyc}}$)}
\label{sec:total_oscillation_time}

As established in \hyperref[app:energy_and_chronology]{Appendix~\ref*{app:energy_and_chronology}}, TFT adopts the FCC site-percolation threshold
\begin{equation}
    p_{\mathrm{c}}
    \approx
    0.199
\end{equation}
under an equal-capacity coarse-graining of the interstitial network.

With
\begin{equation}
    \beta_{\mathrm{IC}}
    \approx
    0.2595,
\end{equation}
the corresponding aggregate live-boundary mass fraction at global percolation is
\begin{equation}
    \beta_{\mathrm{rup}}
    =
    p_{\mathrm{c}}\beta_{\mathrm{IC}}
    \approx
    0.0516.
\end{equation}

The Global Rupture Trigger Mass is therefore
\begin{equation}
    M_{\mathrm{trg}}
    =
    0.0516
    \left(
        1.6707\times10^{79}
        \text{ kg}
    \right)
    \approx
    8.6208\times10^{77}
    \text{ kg}.
\end{equation}

This $5.16\%$ quantity is the aggregate live-boundary mass fraction of the percolating State II network. It is not a $5.16\%$ physical filling of the reservoir by SC-core volume.

The total oscillation time follows from
\begin{equation}
    M_{\mathrm{trg}}
    =
    M_{\mathrm{p}}(0)
    e^{\Gamma_{\mathrm{D}}T_{\mathrm{cyc}}},
\end{equation}
giving
\begin{equation}
    T_{\mathrm{cyc}}
    =
    \frac{1}{\Gamma_{\mathrm{D}}}
    \ln\left(
        \frac{M_{\mathrm{trg}}}
        {M_{\mathrm{p}}(0)}
    \right).
\end{equation}

Because the primordial aggregate must satisfy
\begin{equation}
    M_{\mathrm{p}}(0)
    \geq
    M_{\mathrm{D,G}},
\end{equation}
the minimum-seed normalization,
\begin{equation}
    M_{\mathrm{p}}(0)
    =
    M_{\mathrm{D,G}}
    \approx
    10.033
    \text{ kg},
\end{equation}
defines the maximum possible cycle duration under the present phase-avalanche closure.

For this limiting branch,
\begin{equation}
    8.6208\times10^{77}
    =
    10.033\,
    e^{\Gamma_{\mathrm{D}}T_{\mathrm{cyc}}},
\end{equation}
so
\begin{equation}
    \Gamma_{\mathrm{D}}T_{\mathrm{cyc}}
    =
    \ln\left(
        8.5924\times10^{76}
    \right)
    \approx
    177.15.
\end{equation}

Using
\begin{equation}
    \frac{1}{\Gamma_{\mathrm{D}}}
    \approx
    13.387
    \text{ billion years},
\end{equation}
gives
\begin{equation}
    T_{\mathrm{cyc}}^{(\mathrm{max})}
    \approx
    177.15
    \times
    13.387
    \text{ billion years}
    \approx
    2.3715
    \text{ trillion years}.
\end{equation}

A larger primordial aggregate $M_{\mathrm{p}}(0)$ shortens the absolute elapsed cycle time logarithmically but does not alter the rupture mass or the exponential Global Phase Avalanche law.

\subsection{The Current Cosmic Epoch ($t_{\mathrm{cur}}$)}
\label{sec:current_cosmic_epoch}

To estimate the elapsed time since the last Catastrophic Decompression, we evaluate the Global Phase Avalanche against the present aggregate State II boundary mass,
\begin{equation}
    M_{\mathrm{p}}(t_{\mathrm{cur}})
    \equiv
    M_{\mathrm{cur}}.
\end{equation}

Adopting the present order-of-magnitude Black Hole inventory over the reference observable domain as
\begin{equation}
    M_{\mathrm{obs}}
    \approx
    1.000\times10^{50}
    \text{ kg}
\end{equation}
and using the Black Hole mass-function census as the observational baseline \cite{Sicilia_2022}, TFT extrapolates this inventory across the full Structural Reservoir under the statistical-homogeneity closure developed in \hyperref[app:energy_and_chronology]{Appendix~\ref*{app:energy_and_chronology}}.

Using the structural volumetric scaling factor,
\begin{equation}
    M_{\mathrm{cur}}
    =
    \left(
        1.000\times10^{50}
        \text{ kg}
    \right)
    \left(
        9.7957\times10^{25}
    \right)
    =
    9.7957\times10^{75}
    \text{ kg}.
\end{equation}

The current elapsed phase time is therefore
\begin{equation}
    t_{\mathrm{cur}}
    =
    \frac{1}{\Gamma_{\mathrm{D}}}
    \ln\left(
        \frac{M_{\mathrm{cur}}}
        {M_{\mathrm{p}}(0)}
    \right).
\end{equation}

For the minimum-seed normalization,
\begin{equation}
    9.7957\times10^{75}
    =
    10.033\,
    e^{\Gamma_{\mathrm{D}}t_{\mathrm{cur}}},
\end{equation}
which gives
\begin{equation}
    \Gamma_{\mathrm{D}}t_{\mathrm{cur}}
    \approx
    172.67,
\end{equation}
and
\begin{equation}
    t_{\mathrm{cur}}^{(\mathrm{max})}
    \approx
    172.67
    \times
    13.387
    \text{ billion years}
    \approx
    2.3116
    \text{ trillion years}.
\end{equation}

\vspace{0.5em}
\noindent\textbf{Conclusion:}\\*
Under the minimum-seed normalization, the elapsed phase time satisfies
\begin{equation}
    \frac{t_{\mathrm{cur}}^{(\mathrm{max})}}
    {T_{\mathrm{cyc}}^{(\mathrm{max})}}
    \approx
    0.9747.
\end{equation}

Thus, this limiting chronology places the present epoch at approximately $97.5\%$ of the maximum possible cycle duration. This is a statement about logarithmic elapsed phase time, not about the fraction of the rupture mass already accumulated. The numerical current-age estimate remains conditional on the adopted Black Hole inventory, its extrapolation across the Structural Reservoir, and the single-rate phase-avalanche closure.

\subsection{The Illusion of Genesis (The Cosmic Expiration Timer)}
\label{sec:illusion_of_genesis}

The approximately $13.8$ billion-year age inferred in standard $\Lambda$CDM and the approximately $13.387$ billion-year TFT relaxation time are physically distinct quantities. Standard cosmology obtains its cosmic age from the integrated expansion history, whereas TFT identifies
\begin{equation}
    \frac{1}{\Gamma_{\mathrm{D}}}
\end{equation}
as the characteristic relaxation time of the active vacuum, empirically calibrated by the local Hubble slope \cite{Riess_2022_SH0ES}.

The differential time remaining before Catastrophic Decompression is
\begin{equation}
    \Delta t
    =
    T_{\mathrm{cyc}}
    -
    t_{\mathrm{cur}}.
\end{equation}

Because
\begin{equation}
    M_{\mathrm{trg}}
    =
    M_{\mathrm{p}}(0)
    e^{\Gamma_{\mathrm{D}}T_{\mathrm{cyc}}},
\end{equation}
and
\begin{equation}
    M_{\mathrm{cur}}
    =
    M_{\mathrm{p}}(0)
    e^{\Gamma_{\mathrm{D}}t_{\mathrm{cur}}},
\end{equation}
their ratio eliminates the unknown primordial aggregate exactly:
\begin{equation}
    \frac{M_{\mathrm{trg}}}
    {M_{\mathrm{cur}}}
    =
    e^{\Gamma_{\mathrm{D}}\Delta t}.
\end{equation}

Therefore,
\begin{equation}
    \Delta t
    =
    \frac{1}{\Gamma_{\mathrm{D}}}
    \ln\left(
        \frac{M_{\mathrm{trg}}}
        {M_{\mathrm{cur}}}
    \right).
\end{equation}

Using the adopted global aggregate mass,
\begin{equation}
    \frac{M_{\mathrm{trg}}}
    {M_{\mathrm{cur}}}
    \approx
    88.006,
\end{equation}
so
\begin{equation}
    \Delta t
    \approx
    4.4774
    \frac{1}{\Gamma_{\mathrm{D}}}.
\end{equation}

With the present empirical value of $\Gamma_{\mathrm{D}}$,
\begin{equation}
    \Delta t
    \approx
    59.94
    \text{ billion years}.
\end{equation}

The cancellation of $M_{\mathrm{p}}(0)$ means that this differential clock is independent of the unknown primordial seed normalization. It is not independent of all model assumptions: the numerical value remains conditional on the Global Phase Avalanche closure, the percolation-derived rupture mass, and especially the extrapolated present aggregate mass $M_{\mathrm{cur}}$.

Within those closures, TFT therefore interprets the locally calibrated relaxation rate not as a cosmic birth date but as the ticking rate of a much longer macroscopic structural oscillator.

\subsection{Kinematic Cycle Reach and the Optical Assimilation Boundary}
\label{sec:causal_distance}

The thermodynamic chronology defines a second macroscopic length scale distinct from both $R_{\mathrm{U}}$ and $R_{\mathrm{Res}}$.

For a transverse excitation propagating at $c_{\mathrm{E}}$, the complete structural-cycle duration corresponds to the kinematic cycle length
\begin{equation}
    L_{\mathrm{max}}
    \equiv
    c_{\mathrm{E}}T_{\mathrm{cyc}}.
\end{equation}

For the minimum-seed, maximum-cycle branch,
\begin{equation}
    L_{\mathrm{max}}
    \approx
    2.2436\times10^{28}
    \text{ m}.
\end{equation}

This is the greatest propagation-equivalent distance available to a $c_{\mathrm{E}}$ transverse mode over the maximum structural-cycle duration. It is not an absolute causal boundary of the substrate.

The distinction is essential because the physical Structural Reservoir at rupture is much larger,
\begin{equation}
    R_{\mathrm{rup}}
    \approx
    5.7361\times10^{34}
    \text{ m},
\end{equation}
so
\begin{equation}
    L_{\mathrm{max}}
    \ll
    R_{\mathrm{rup}}.
\end{equation}

Consequently, ordinary $c_{\mathrm{E}}$ transverse propagation cannot traverse the entire Structural Reservoir within one complete cycle. This establishes a finite transverse communication reach per structural epoch, but it does not imply absolute causal disconnection of the underlying substrate or eliminate the distinct longitudinal Hard-Floor structural channel.

Moreover, $L_{\mathrm{max}}$ is generally not the physical coherent-photon extinction distance. Coherent optical propagation is instead limited by the frequency-dependent CMB-reservoir assimilation condition. With
\begin{equation}
    \omega(x)
    =
    \omega_{\mathrm{emit}}
    e^{-x/R_{\mathrm{U}}},
\end{equation}
the assimilation distance satisfies
\begin{equation}
    \omega(x_{\mathrm{assim}})
    =
    \omega_{\mathrm{CMB}}^{(\mathrm{eff})},
\end{equation}
and therefore
\begin{equation}
    x_{\mathrm{assim}}
    =
    R_{\mathrm{U}}
    \ln\left(
        \frac{\omega_{\mathrm{emit}}}
        {\omega_{\mathrm{CMB}}^{(\mathrm{eff})}}
    \right).
\end{equation}

Thus, the CMB noise-floor boundary is frequency dependent and may occur well before $L_{\mathrm{max}}$. The three scales have distinct meanings:
\begin{equation}
\begin{array}{rcl}
    R_{\mathrm{U}}
    &:&
    \text{one-$e$-fold coherent-relaxation scale},
    \\[3pt]
    x_{\mathrm{assim}}
    &:&
    \text{frequency-dependent coherent assimilation scale},
    \\[3pt]
    L_{\mathrm{max}}
    &:&
    \text{kinematic length corresponding to the structural cycle}.
\end{array}
\end{equation}

Their common dependence on $\Gamma_{\mathrm{D}}$ links the optical and thermodynamic sectors without identifying these physically distinct boundaries.

\subsection{Macroscopic Clock Symmetry and Characteristic Transverse-Reach Domains}

The limited transverse reach of a single structural epoch is supported by two complementary macroscopic relations.

\vspace{0.5em}
\noindent\textbf{1. The Macroscopic Phase Symmetry (Optical and Thermodynamic Clock Equivalence)}\\*

As derived in \hyperref[app:energy_and_chronology]{Appendix~\ref*{app:energy_and_chronology}}, the shared appearance of $\Gamma_{\mathrm{D}}$ in coherent optical relaxation and in the Global Phase Avalanche produces a reciprocal clock identity.

Because physical CMB-reservoir assimilation may occur before $L_{\mathrm{max}}$, the optical exponential is formally continued to the structural endpoint:
\begin{equation}
    E_{\mathrm{obs}}^{(\mathrm{formal})}(L_{\mathrm{max}})
    =
    E_{\mathrm{emit}}
    \exp\left(
        -\frac{\Gamma_{\mathrm{D}}L_{\mathrm{max}}}
        {c_{\mathrm{E}}}
    \right).
\end{equation}

Since
\begin{equation}
    L_{\mathrm{max}}
    =
    c_{\mathrm{E}}T_{\mathrm{cyc}},
\end{equation}
we obtain
\begin{equation}
    \frac{
        E_{\mathrm{obs}}^{(\mathrm{formal})}(L_{\mathrm{max}})
    }{
        E_{\mathrm{emit}}
    }
    =
    e^{-\Gamma_{\mathrm{D}}T_{\mathrm{cyc}}}.
\end{equation}

At the same structural endpoint, the Global Phase Avalanche gives
\begin{equation}
    \frac{M_{\mathrm{trg}}}
    {M_{\mathrm{p}}(0)}
    =
    e^{+\Gamma_{\mathrm{D}}T_{\mathrm{cyc}}}.
\end{equation}

Therefore,
\begin{equation}
    \frac{
        E_{\mathrm{obs}}^{(\mathrm{formal})}(L_{\mathrm{max}})
    }{
        E_{\mathrm{emit}}
    }
    =
    \frac{
        M_{\mathrm{p}}(0)
    }{
        M_{\mathrm{trg}}
    }.
\end{equation}

For the minimum-seed normalization,
\begin{equation}
    M_{\mathrm{p}}(0)
    =
    M_{\mathrm{D,G}},
\end{equation}
this becomes
\begin{equation}
    \frac{
        E_{\mathrm{obs}}^{(\mathrm{formal})}(L_{\mathrm{max}})
    }{
        E_{\mathrm{emit}}
    }
    =
    \frac{
        M_{\mathrm{D,G}}
    }{
        M_{\mathrm{trg}}
    }.
\end{equation}

If the emitted excitation is normalized to
\begin{equation}
    E_{\mathrm{emit}}
    =
    E_{\mathrm{max}}
    \approx
    13.331
    \text{ J},
\end{equation}
then
\begin{equation}
    E_{\mathrm{obs}}^{(\mathrm{formal})}(L_{\mathrm{max}})
    =
    E_{\mathrm{max}}
    \left(
        \frac{M_{\mathrm{D,G}}}
        {M_{\mathrm{trg}}}
    \right)
    \approx
    1.5515\times10^{-76}
    \text{ J}.
\end{equation}

This quantity is a formal continuation of the optical relaxation clock, not necessarily the energy of an individual coherent photon physically surviving to $L_{\mathrm{max}}$. The reciprocal relation likewise does not imply direct energy transfer between the optical carrier and the State II network. It expresses the opposite exponential evolution of two distinct processes under the same single-rate macroscopic closure.

\vspace{0.5em}
\noindent\textbf{2. The Number of Characteristic Transverse-Reach Domains ($N_{\mathrm{dom}}$)}\\*

The large separation between $R_{\mathrm{rup}}$ and $L_{\mathrm{max}}$ may be expressed as a characteristic volume ratio. Comparing the spherical rupture volume with a sphere of radius $L_{\mathrm{max}}$ gives
\begin{equation}
    N_{\mathrm{dom}}
    \equiv
    \frac{V_{\mathrm{rup}}}
    {V_{L_{\mathrm{max}}}}
    =
    \left(
        \frac{R_{\mathrm{rup}}}
        {L_{\mathrm{max}}}
    \right)^3.
\end{equation}

Substituting the maximum-cycle values,
\begin{equation}
    N_{\mathrm{dom}}
    =
    \left(
        \frac{
            5.7361\times10^{34}
        }{
            2.2436\times10^{28}
        }
    \right)^3
    \approx
    \left(
        2.5567\times10^6
    \right)^3,
\end{equation}
which gives
\begin{equation}
    N_{\mathrm{dom}}
    \approx
    1.67\times10^{19}.
\end{equation}

Thus, the Structural Reservoir contains the volume equivalent of approximately $16.7$ quintillion characteristic transverse-causal domains.

This expression is a volume-ratio measure; it is not a claim that equal $L_{\mathrm{max}}$ spheres geometrically tessellate the reservoir.

Its causal interpretation is nevertheless physical. Regions separated by more than the available $c_{\mathrm{E}}$ transverse reach cannot establish direct transverse communication during the same State I cycle.

They remain, however, components of one contiguous mechanical substrate.

The TFT universe is therefore not globally disconnected in the structural sense. Its ordinary transverse causal structure is local while its mechanical constitution is global.

That distinction permits a Structural Reservoir of order $10^{35}$ m to undergo a collective phase failure even though ordinary electromagnetic information cannot traverse the reservoir before the cycle resets.

